\section{本章纲要：征服重组的是区域关系}

蒙古扩张不能只按胜者的行军路线叙述。从西夏的绿洲与金的华北，到南宋的四川、江汉、江南和海上，战争面对的是不同的道路、城防、税源、技术与地方政治。蒙古诸部的联盟、军事组织和吸纳能力固然重要，但投附者、抵抗者、工匠、商人、官员和被迁徙家庭同样决定征服能否从一次战役变成区域接收。

本章追问旧政权留下的城市、文书、军户、港口和交通网络怎样被借用、改造或毁坏。西夏、金和南宋的终结各有长期的社会时间，不能由一个亡国年份抹平。官修史、多语叙事、地方文书、碑刻与考古材料将被并置使用：它们能够确认若干节点，也会保留战损、人数和动机的争议，避免让后来的帝国秩序替所有地区说明征服。
